\documentclass[12pt]{article}
\usepackage[margin=1in]{geometry}
\usepackage{setspace}
\onehalfspacing{}
\usepackage[dvipsnames,table,xcdraw]{xcolor} % colors

% Start of preamble
%==========================================================================================%
% Required to support mathematical unicode
\usepackage[warnunknown, fasterrors, mathletters]{ucs}
\usepackage[utf8x]{inputenc}

% Standard mathematical typesetting packages
\usepackage{amsmath,amssymb,amscd,amsthm,amsxtra, pxfonts}
\usepackage{mathtools,mathrsfs,xparse}

% Symbol and utility packages
\usepackage{cancel, textcomp}
\usepackage[mathscr]{euscript}
\usepackage[nointegrals]{wasysym}
\usepackage{apacite}

% Extras
\usepackage{physics}  % Lots of useful shortcuts and macros
\usepackage{tikz-cd}  % For drawing commutative diagrams easily
\usepackage{microtype}  % Minature font tweaks
%\usepackage{pgfplots} % plots

\usepackage{enumitem}
\usepackage{titling}

\usepackage{graphicx}

%\usepackage{quiver}

% Fancy theorems due to @intuitively on discord
\usepackage{mdframed}
\newmdtheoremenv[
backgroundcolor=NavyBlue!30,
linewidth=2pt,
linecolor=NavyBlue,
topline=false,
bottomline=false,
rightline=false,
innertopmargin=10pt,
innerbottommargin=10pt,
innerrightmargin=10pt,
innerleftmargin=10pt,
skipabove=\baselineskip,
skipbelow=\baselineskip]{mytheorem}{Theorem}

\newenvironment{theorem}{\begin{mytheorem}}{\end{mytheorem}}

\newtheorem{corollary}{Corollary}
\newtheorem{lemma}{Lemma}

\newtheoremstyle{definitionstyle}
{\topsep}%
{\topsep}%
{}%
{}%
{\bfseries}%
{.}%
{.5em}%
{}%
\theoremstyle{definitionstyle}
\newmdtheoremenv[
backgroundcolor=Violet!30,
linewidth=2pt,
linecolor=Violet,
topline=false,
bottomline=false,
rightline=false,
innertopmargin=10pt,
innerbottommargin=10pt,
innerrightmargin=10pt,
innerleftmargin=10pt,
skipabove=\baselineskip,
skipbelow=\baselineskip,
]{mydef}{Definition}
\newenvironment{definition}{\begin{mydef}}{\end{mydef}}

\newtheorem*{remark}{Remark}

\newtheorem*{example}{Example}
\newtheorem*{claim}{Claim}

% Common shortcuts
\def\mbb#1{\mathbb{#1}}
\def\mfk#1{\mathfrak{#1}}

\def\bN{\mbb{N}}
\def\C{\mbb{C}}
\def\R{\mbb{R}}
\def\bQ{\mbb{Q}}
\def\bZ{\mbb{Z}}
\def\cph{\varphi}
\renewcommand{\th}{\theta}
\def\ve{\varepsilon}
\newcommand{\mg}[1]{\| #1 \|}

% Often helpful macros
\newcommand{\floor}[1]{\left\lfloor#1\right\rfloor}
\newcommand{\ceil}[1]{\left\lceil#1\right\rceil}
\renewcommand{\qed}{\hfill\qedsymbol}
\renewcommand{\ip}[1]{\langle#1\rangle}
\newcommand{\seq}[2]{\qty(#1_#2)_{#2=1}^{\infty}}

\newcommand{\SET}[1]{\Set{\mskip-\medmuskip #1 \mskip-\medmuskip}}

% End of preamble
%==========================================================================================%

% Start of commands specific to this file
%==========================================================================================%

\usepackage{braket}
\usepackage{dsfont}
\newcommand{\Z}{\mbb Z}
\newcommand{\gen}[1]{\left\langle #1 \right\rangle}
\newcommand{\nsg}{\trianglelefteq}
\newcommand{\F}{\mbb F}
\newcommand{\Aut}{\mathrm{Aut}}
\newcommand{\sepdeg}[1]{[#1]_{\mathrm{sep}}}
\newcommand{\Q}{\mbb Q}
\newcommand{\Gal}{\mathrm{Gal}\qty}
\newcommand{\id}{\mathrm{id}}
\newcommand{\Hom}{\mathrm{Hom}_R}
\newcommand{\1}{\mathds 1}
\newcommand{\N}{\mathbb N}
\renewcommand{\P}{\mathbb P \qty}
\newcommand{\E}{\mathbb E \qty}
\newcommand{\Var}{\mathrm{Var}}
\everymath{\displaystyle}
\newcommand{\argmax}{\mathrm{argmax}}

%==========================================================================================%
% End of commands specific to this file

\title{Math Template}
\date{\today}
\author{Rohan Mukherjee}

\begin{document}
    \maketitle
    \begin{enumerate}
        \item I give two proofs for this first assertion. The first is by induction. The base case, that $\P(S_1 = X_1 \leq x) = x$ for $0 \leq x \leq 1$ is true by definition. Then suppose that for $0 \leq x \leq 1$, $\P(S_n \leq x) = \frac{x^n}{n!}$. Then for $0 \leq x \leq 1$, we have that:
        \begin{align*}
            \P(S_{n+1} \leq x) &= \int_0^x \int_{-\infty}^\infty p_{X_{n+1}}(x-y) p_{S_n}(y) \dd y \dd x \\
            &= \int_0^x \int_0^x \frac{s^{n-1}}{(n-1)!} \dd x \dd y \\
            &= \int_0^x \frac{s^n}{n!} \dd x \\
            &= \frac{x^{n+1}}{(n+1)!}
        \end{align*}
        In the middle we used that $p_{S_n}(x) = \dv{x} \P(S_n \leq x) = \frac{x^{n-1}}{(n-1)!}$. This completes the induction. 

        Plugging in, we get $\P(S_n \leq n) = \frac{1}{n!}$. 

        The second proof is by taking cross-sections. First, we know that the area of an isocoles right triangle with side lengths $x$ has area $\frac{x^2}{2}$. The horizontal cross sections of the tetrahedron defined by $x+y+z \leq 1$ are isocoles right triangles. By summing all of these, we get that the area of this tetrahedron is:
        \begin{align*}
            \int_0^1 \int_0^{1-x} \int_0^{1-x-y} \dd z \dd y \dd x &= \int_0^1 \int_0^{1-x} (1-x-y) \dd y \dd x \\
            &= \int_0^1 \frac{(1-x)^2}{2} \dd x \\
            &= \frac{1}{6}
        \end{align*}
        Let the volume of the $n$-simplex be $V_n$. Taking a cross-section of the $n+1$-simplex defined by $x_1 + \cdots + x_{n+1} \leq 1$ will yield an $n$-simplex. If one of the side lengths of this $n$-simplex is $x$, then by similarity with the bottom $n$-simplex who has side lengths of 1, the area of this $n$-simplex is $x^nV_n$. By summing all these up, we get that the volume of the $n+1$-simplex is:
        \begin{align*}
            \int_0^1 x^n V_n \dd x = \frac{V_n}{n+1}
        \end{align*}
        Induction yields that $V_n = \frac{1}{n!}$.

        The third proof is a little bit of linear algebra, and my favorite one. Consider $S = \SET{0 \leq x_1 \leq x_2 \leq \cdots \leq x_n \leq 1}$. Then for any permutation $\pi \in S_n$, the set $\pi(S) = \SET{0 \leq x_{\pi(1)} \leq x_{\pi(2)} \leq \cdots \leq x_{\pi(n)} \leq 1}$ has the same volume as $S$. Since the disjoint union of these sets is $[0,1]^n$, we get that the volume of $S$ is just $\frac{1}{n!}$. 

        This relates to the original simplex in the following way. Let $A$ be the linear map defined by $A(x_1, \ldots, x_n) = (x_1, x_2 - x_1, \ldots, x_n - x_{n-1})$. A simple exercise shows that $A(S) = \SET{x_1 + \cdots + x_n \leq 1}$. Then by high-dimensional geometry, the volume of $A(S)$ is $\det(A) \cdot S$. The matrix of $A$ is of the following form:
        \begin{align*}
            \begin{pmatrix}
                1 & 0 & 0 & \cdots & 0 \\
                -1 & 1 & 0 & \cdots & 0 \\
                0 & -1 & 1 & \cdots & 0 \\
                \vdots & \vdots & \vdots & \ddots & \vdots \\
                0 & 0 & 0 & \cdots & 1
            \end{pmatrix}
        \end{align*}
        Starting from the identity matrix, by subtracting the $i$th row from the $i+1$th row, which preserves the determinant, we get that $\det(A) = 1$. Thus the volume of $\SET{x_1 + \cdots + x_n \leq 1}$ is $\frac{1}{n!}$.

        Now knowing that $\P(T \geq n) = \frac{1}{n!}$, we can compute the expected value of $T$ as:
        \begin{align*}
            \E[T] = \sum_{n=1}^\infty \P(T \geq n) = \sum_{n=1}^\infty \frac{1}{n!} = e
        \end{align*}

        Notice that $S_T = \sum_{n=1}^T X_n = \sum_{n=1}^\infty X_n \1_{T \geq n}$. Then as $S_T \geq 0$ and Fubini-Tonelli theorem, we can switch the order of summation and integration to get:
        \begin{align*}
            \E[\sum_{n=1}^\infty X_n \1_{T \geq n}] = \sum_{n=1}^\infty \E[X_n] \P(T \geq n) \\
            &= \sum_{n=1}^\infty \frac 12 \frac{1}{n!} = \frac{e}{2}
        \end{align*}
    \end{enumerate}
\end{document}